% Created 2026-06-13 Sat 14:03
% Intended LaTeX compiler: xelatex
\documentclass[11pt]{article}
\usepackage{graphicx}
\usepackage{longtable}
\usepackage{wrapfig}
\usepackage{rotating}
\usepackage[normalem]{ulem}
\usepackage{capt-of}
\usepackage{hyperref}
\usepackage[a4paper,margin=0.75in,top=0.85in]{geometry}
\usepackage{amsmath,amssymb,mathtools}
\usepackage{microtype}
\setlength{\parskip}{0.8em}
\setlength{\parindent}{0pt}
\usepackage{mathpazo}
\usepackage{titlesec}
\titleformat{\section}{\large\bfseries}{}{0em}{}[\titlerule]
\titleformat{\subsection}{\normalsize\bfseries}{}{0em}{}
\usepackage{enumitem}
\setlist{noitemsep,topsep=4pt}
\usepackage{titling}
\pretitle{\begin{center}\large\bfseries}
\posttitle{\end{center}}
\author{Tyrone}
\date{\today}
\title{Rules of Inference}
\hypersetup{
 pdfauthor={Tyrone},
 pdftitle={Rules of Inference},
 pdfkeywords={},
 pdfsubject={},
 pdfcreator={Emacs 30.2 (Org mode 9.7.11)}, 
 pdflang={English}}
\begin{document}

\maketitle
\begin{enumerate}
\item \textbf{Modus Ponens} (Law of Detachment)

If a conditional statement is true and its condition is satisfied, the conclusion must be true.
\[
   \text{If } p \rightarrow q \text{ and } p, \text{ then } q
   \]

\item \textbf{Modus Tollens} (Law of Contrapositive)

If a conditional statement is true, and its consequent is false, then its antecedent must also be false.
\[
   \text{If } p \rightarrow q \text{ and } \neg q, \text{ then } \neg p
   \]

\item \textbf{Hypothetical Syllogism}

If two conditional statements are true, where the consequent of the first is the antecedent of the second, then a third conditional statement combining the antecedent of the first and the consequent of the second is also true.
\[
   \text{If } p \rightarrow q \text{ and } q \rightarrow r, \text{ then } p \rightarrow r
   \]

\item \textbf{Disjunctive Syllogism}
If a disjunction is true, and of its parts is false, then the other part must be  true.

\[
   \text{If } p \lor q \text{ and } \neg p \text{, then } q
   \]

\item \textbf{Conjunction}

If two statements are true, then their conjunction is also true.
\[
   \text{If } p \text{ and } q, \text{ then } p \land q
   \]

\item \textbf{Simplification}

If a conjunction is true, then each of its conjuncts are also true.
\[
   \text{If } p \land q, \text{ then } p
   \]

\item \textbf{Addition}

If a statement is true, then the disjunction of that statement with any other statement is also true.
\[
   \text{If } p, \text{ then } p \lor q
   \]

\item \textbf{Absorption}

If a conditional statement is true, then the antecedent implies a conjunction of itself and the consequent.
\[
   \text{If } p \rightarrow q, \text{ then } p \rightarrow (p \land q)
   \]

\item \textbf{Resolution}

If two disjunctions are true, and one containing proposition p while the other contains its negation \(\neg p\), then the disjunction of the remaining parts is true.
\[
   \text{If } p \lor q \text{ and } \neg p \lor r, \text{ then } q \lor r
   \]
\end{enumerate}
\end{document}
